\documentclass[11pt]{article}
\usepackage[margin=1in]{geometry}
\usepackage{amsmath,amssymb,amsthm,mathtools}
\usepackage[hidelinks]{hyperref}
\usepackage{enumitem}
\newtheorem{lemma}{Lemma}
\newtheorem{proposition}{Proposition}
\newtheorem{corollary}{Corollary}
\newcommand{\E}{\mathbb E}
\newcommand{\cE}{\mathcal E}
\newcommand{\cP}{\mathcal P}
\title{Erd\H{o}s Problem 625: Result Extensions and Proof Simplifications}
\author{Research note for the candidate manuscript}
\date{24 July 2026}
\begin{document}
\maketitle

\paragraph{Status.}
This note records exact deductions, finite rational certificates, and
architectural alternatives obtained by re-examining the candidate manuscript
at repository main commit
\begin{center}
\small\ttfamily cda78922ea6c87bfc81f9bf693374dd045dac624
\end{center}
It is not external peer review, priority verification, or a completed Lean
proof of the final target.  Put $q=\ln2$ and $N=\ln n$.

\section{Exact simplification of the large-residual attachment bound}

\begin{lemma}[An even completion of residual edges is unique]
Let $M$ be a matching and $R$ any other finite edge set.  Put
\[
 \cE(M,R)=\{F\subseteq M\cup R:\deg_F(v)\equiv0\pmod2\text{ for every }v\}.
\]
Then
\[
 \rho:\cE(M,R)\to\cP(R\setminus M),\qquad \rho(F)=F\setminus M
 \tag{1.1}
\]
is injective.
\end{lemma}

\begin{proof}
If $\rho(F)=\rho(F')$, then $F\mathbin\triangle F'\subseteq M$.  The symmetric
difference is even.  A nonempty subset of a matching has degree one at every
incident vertex, so it is not even.  Thus $F\mathbin\triangle F'=\varnothing$.
\end{proof}

\begin{corollary}[Weighted product bound]
For nonnegative residual weights $(q_e)$, extended by zero on $M$,
\[
 \sum_{F\in\cE(M,R)}\prod_{e\in F\setminus M}q_e
 \le\prod_{e\in R\setminus M}(1+q_e)
 \le\exp\left(\sum_e q_e\right).
 \tag{1.2}
\]
\end{corollary}

Equation (9.12) therefore gives
\[
 \mathcal A(M,j)\le\exp\left(\Lambda_0+\sum_eq_e\right).
 \tag{1.3}
\]
Since
\[
 \sum_eq_e=\frac12\sum_{a,b}\theta_{ab}^2+\Lambda_0
 \tag{1.4}
\]
and
\[
 \sum_{a,b}\theta_{ab}^2
 =\frac{e^2}{m_0^2}
   \left(\sum_ad_a^2\right)\left(\sum_b(d'_b)^2\right)
 \le e^2U^2,
 \tag{1.5}
\]
while $\Lambda_0\le CU^4/m_0$, we obtain
\[
 \boxed{\mathcal A(M,j)
 \le\exp\left\{C\left(U^2+\frac{U^4}{m_0}\right)\right\}.}
 \tag{1.6}
\]
For $m_0\ge n/N^6$ and $U=O(N)$ this is $\exp(CN^2)$, stronger than the
current $\exp(CN^8)$ bound.  It removes the simple-cycle decomposition,
residual and mixed walk enumeration, $\tau$, and equations (9.15)--(9.18).

\section{A stronger central partial-diagonal rate}

\begin{lemma}
Under the hypotheses of Lemma 7.1A,
\[
 \Phi_T(z)\le-\frac{1-R}{100}\qquad(1/64\le R\le1).
 \tag{2.1}
\]
\end{lemma}

\begin{proof}
On $[1/64,47/100]$ use
\[
 \Phi_T(z)\le R\ln R+(5q/2-1)R.
\]
The convex function
\[
 f(R)=R\ln R+(5q/2-1)R+\frac{1-R}{100}
\]
has $f(1/64)<-0.0436$ and $f(47/100)<-0.0051$.  On
$[47/100,1]$ use
\[
 \Phi_T(z)\le R\ln R+(1-q/2)(1-R).
\]
The convex function
\[
 h(R)=R\ln R+(1-q/2+1/100)(1-R)
\]
has $h(47/100)<-0.0032$ and $h(1)=0$.
\end{proof}

\section{Carry the phase-resolved displacement to the theorem}

Put
\[
 A(\delta)=q-D_4(\delta),\qquad
 A_*=\min_{0\le\delta\le1}A(\delta).
 \tag{3.1}
\]
Continuity and Lemma 5.1 give $A_*>\gamma_4$, where
$\gamma_4=\ln(200/153)$.  Equation (5.11), the midpoint construction, and the
$o(n/N^3)$ amplification loss give
\[
 \boxed{\chi(G_n)-\zeta(G_n)
 \ge\left(\frac{q^2}{8}A(\delta_n)-o(1)\right)\frac{n}{N^3}}
 \tag{3.2}
\]
with high probability.  Thus the fixed displayed constant can be increased
from $q^2\gamma_4/32$ to $q^2\gamma_4/8$ without changing the witness or
second-moment argument.

\section{A stronger four-support entropy certificate}

\begin{proposition}
For every phase,
\[
 \frac{12}{5}q<\lambda_4<\frac{21}{5}q.
 \tag{4.1}
\]
Moreover
\[
 L(12q/5)<3/10,\quad H(3q)<1/50,
 \quad L(3q)<3/25,\quad H(21q/5)<1/5.
 \tag{4.2}
\]
Consequently
\[
 D_4(\delta)<\ln(33/25),\qquad q-D_4(\delta)>\ln(50/33).
 \tag{4.3}
\]
\end{proposition}

Put $x=2^{1/10}$.  The rational intervals
\[
 0.693147<q<0.693148,
 \qquad 1.071773<x<1.071774
 \tag{4.4}
\]
reduce the new comparisons to integer powers.  For example,
\[
 10(1+x^{29}+x^{48})
 <3(x^{57}+x^{56}+x^{45}+x^{24})
 \tag{4.5}
\]
gives $L(12q/5)<3/10$, while
\[
 H(21q/5)\le
 \frac{x^{72}}{(1-x^{-23})(x^{64}+x^{81}+x^{88}+x^{85})}<\frac15.
 \tag{4.6}
\]
The standard-library script \texttt{entropy\_certificate\_upgrade.py} checks
these rational statements and the two mean brackets exactly.

Combining this certificate with (3.2) gives
\[
 \boxed{c_{\rm improved}
 =\frac{(\ln2)^2}{8}\ln\frac{50}{33}
 =0.0249544559\ldots,}
 \tag{4.7}
\]
more than six times the current displayed coefficient.

\section{Complement-symmetric strengthening}

Since $\zeta(\overline G)=\zeta(G)$ and $G(n,1/2)$ is complement invariant,
intersecting the theorem for $G_n$ and $\overline G_n$ gives
\[
 \boxed{\min\{\chi(G_n),\chi(\overline G_n)\}-\zeta(G_n)
 \ge c\frac{n}{(\ln n)^3}}
 \tag{5.1}
\]
with high probability.

\section{Fixed-density criticality away from one half}

For every graph,
\[
 \zeta(G)\le\min\{\chi(G),\chi(\overline G)\}.
\]
For fixed $p\in(0,1)$,
\[
 \chi(G(n,p))
 =\left(\frac12\ln\frac1{1-p}+o(1)\right)\frac{n}{\ln n}.
\]
Therefore
\[
 \boxed{\max\{\chi(G),\chi(\overline G)\}-\zeta(G)
 \ge\left(\frac12\left|\ln\frac{p}{1-p}\right|+o(1)\right)
 \frac{n}{\ln n}.}
 \tag{6.1}
\]
For $p>1/2$ this is a direct lower bound on $\chi(G)-\zeta(G)$; for $p<1/2$
it applies to the complementary chromatic gap.  Thus $p=1/2$ is the unique
fixed-density point where the first-order comparison cancels.

\section{A standard upper bound}

If $h(G)=\max\{\alpha(G),\omega(G)\}$, then $\zeta(G)\ge n/h(G)$.  The standard
second-order estimates at $p=1/2$ give
\[
 h(G)=2\log_2n-2\log_2\log_2n+O(1)
\]
and
\[
 \chi(G)=\frac{n}{2\log_2n-2\log_2\log_2n+O(1)}.
\]
Hence
\[
 \boxed{\chi(G)-\zeta(G)=O\left(\frac{n}{(\ln n)^2}\right)}
 \tag{7.1}
\]
with high probability.

\section{Architectural alternatives}

Numerical optimization shows that the three-size support $\{2,3,5\}$ still has
uniform positive signed advantage, approximately $0.0921449643$.  It preserves
the unimodular correction coordinates $2,3$, the central-rate endpoints $2,5$,
and maximum type displacement three, while reducing the transportation table
from $4\times4$ to $3\times3$.  This is a diagnostic, not yet a proved
replacement route.

The midpoint can also be replaced by
\[
 k_\theta=\left\lceil r_4^{co}+\theta(r_+-r_4^{co})\right\rceil
\]
for fixed $\theta\in(0,1)$.  The retained gap would be
\[
 \left((1-\theta)\frac{q^2}{4}A(\delta)+o(1)\right)\frac{n}{N^3},
\]
while the signed first moment keeps a fixed positive exponential margin.  This
requires a route-by-route replay of every use of $c_Z$ before being promoted to
a theorem.

\section{Recommended order}
\begin{enumerate}
\item Replace the large-residual cycle expansion by the restriction injection.
\item Strengthen the central rate.
\item Carry equation (5.11) directly to the conclusion.
\item Review and insert the exact entropy certificate.
\item Add the complement and fixed-density corollaries.
\item Keep the three-size and non-midpoint variants on experimental branches.
\end{enumerate}
\end{document}
